O  projeto tem como objetivo o desenvolvimento de um sistema de monitoramento capaz de coletar dados de temperatura e umidade do ambiente e disponibilizá-los em tempo real por meio de uma interface web.

O sistema será composto por um sensor de temperatura e umidade, responsável pela aquisição dos dados ambientais, e por um microcontrolador, encarregado de processar essas informações, realizar a comunicação com a interface web e controlar os elementos de saída do sistema, como os LEDs indicadores de estado. A escolha desses componentes será fundamentada em critérios técnicos relevantes, como faixa de medição, precisão, consumo de energia, facilidade de integração e capacidade de comunicação em rede.

Além da aquisição e exibição dos dados, o sistema deverá permitir a interação com o usuário por meio de uma interface web, na qual será possível visualizar as medições em tempo real e alternar entre diferentes modos de operação. Os modos previstos são: desligado, no qual o sistema permanece inativo; automático, em que a coleta e exibição dos dados ocorrem de forma contínua; e ligado, no qual o sistema aguarda ações do usuário. A mudança de estado deverá ser refletida visualmente através de LEDs, proporcionando uma indicação clara e imediata do modo de operação atual.

O desenvolvimento do projeto envolve a integração entre hardware e software, implementando a lógica de funcionamento do sistema e da interface web. Dessa forma, busca-se construir uma solução funcional, organizada e alinhada aos requisitos propostos.

